\section{Windows Investigations}

Windows systems generate some of the most important telemetry in enterprise environments. 
Because adversaries frequently target Windows endpoints, domain controllers, and servers, 
understanding Windows event logs is essential for threat hunting, detection engineering, 
and CTI validation. This chapter introduces the core Windows telemetry sources, their most 
important fields, and practical investigation workflows that support real-world analysis.

% --------------------------------------------------------------------
% SECTION 1: Core Windows Telemetry Sources
% --------------------------------------------------------------------

\subsection*{Core Windows Telemetry Sources}

Windows produces several categories of logs relevant to security investigations. 
Each provides unique visibility into attacker behavior:

\begin{itemize}
    \item \textbf{Security Event Log} — Authentication, authorization, account changes, policy modifications.
    \item \textbf{Sysmon} — High-fidelity process, network, registry, file, and DNS telemetry.
    \item \textbf{PowerShell Operational Log} — Script block logging, module loading, command execution.
    \item \textbf{WMI Activity Log} — WMI persistence, remote execution, system management operations.
    \item \textbf{Task Scheduler Logs} — Scheduled task creation, modification, and execution.
\end{itemize}

These sources form the foundation of Windows investigations.

% --------------------------------------------------------------------
% SECTION 2: Security Event Log Fundamentals
% --------------------------------------------------------------------

\subsection*{Security Event Log Fundamentals}

The Windows Security Event Log contains authentication events, account activity, and policy changes. 
Some of the most important Event IDs include:

\begin{itemize}
    \item \textbf{4624} — Successful logon
    \item \textbf{4625} — Failed logon
    \item \textbf{4672} — Special privileges assigned
    \item \textbf{4720} — User account created
    \item \textbf{4722} — User account enabled
    \item \textbf{4723/4724} — Password reset or change
    \item \textbf{4768/4769} — Kerberos ticket activity
\end{itemize}

These events are essential for identity investigations, brute-force detection, and lateral movement analysis.

\textbf{Example SPL Query}

\begin{lstlisting}
index=wineventlog EventCode=4625
| stats count by AccountName, WorkstationName
| where count > 10
\end{lstlisting}

% --------------------------------------------------------------------
% SECTION 3: Sysmon Fundamentals
% --------------------------------------------------------------------

\subsection*{Sysmon Fundamentals}

Sysmon provides detailed endpoint telemetry that is invaluable for threat hunting. 
Key Event Codes include:

\begin{itemize}
    \item \textbf{1} — Process creation
    \item \textbf{3} — Network connection
    \item \textbf{7} — Image loaded
    \item \textbf{11} — File created
    \item \textbf{13} — Registry modification
    \item \textbf{22} — DNS query
\end{itemize}

Sysmon is particularly useful for detecting execution chains, malware behavior, and suspicious parent-child relationships.

\textbf{Example SPL Query}

\begin{lstlisting}
index=sysmon EventCode=1
| table _time, Computer, User, ParentImage, Image, CommandLine
\end{lstlisting}

% --------------------------------------------------------------------
% SECTION 4: PowerShell Logging
% --------------------------------------------------------------------

\subsection*{PowerShell Logging}

PowerShell is a common tool for both administrators and adversaries. 
The most important PowerShell logs include:

\begin{itemize}
    \item \textbf{4103} — Module logging
    \item \textbf{4104} — Script block logging
    \item \textbf{600} — Engine lifecycle events
\end{itemize}

Script block logging (4104) is especially valuable for detecting encoded or obfuscated commands.

\textbf{Example SPL Query}

\begin{lstlisting}
index=wineventlog EventCode=4104
| search CommandLine="*EncodedCommand*"
\end{lstlisting}

% --------------------------------------------------------------------
% SECTION 5: WMI Investigations
% --------------------------------------------------------------------

\subsection*{WMI Investigations}

WMI is frequently used for remote execution, persistence, and reconnaissance. 
Relevant logs include:

\begin{itemize}
    \item \textbf{Microsoft-Windows-WMI-Activity/Operational}
    \item \textbf{WMI-Activity}
\end{itemize}

\textbf{Example SPL Query}

\begin{lstlisting}
index=wineventlog sourcetype="WMI-Activity"
| stats count by Operation, User
\end{lstlisting}

% --------------------------------------------------------------------
% SECTION 6: Common Investigation Workflows
% --------------------------------------------------------------------

\subsection*{Common Investigation Workflows}

Windows investigations often follow repeatable patterns. 
The following SPL examples demonstrate practical workflows used in threat hunting and detection engineering.

\textbf{1. Process Chain Analysis}

\begin{lstlisting}
index=sysmon EventCode=1
| stats count by ParentImage, Image
| where count < 5
\end{lstlisting}

\textbf{2. Authentication Pattern Analysis}

\begin{lstlisting}
index=wineventlog EventCode=4625
| stats count by AccountName, WorkstationName
\end{lstlisting}

\textbf{3. PowerShell Abuse Detection}

\begin{lstlisting}
index=wineventlog EventCode=4104
| search CommandLine="*EncodedCommand*"
\end{lstlisting}

% --------------------------------------------------------------------
% SECTION 7: Why Windows Investigations Matter
% --------------------------------------------------------------------

\subsection*{Why Windows Investigations Matter}

Windows telemetry provides visibility into authentication, execution, persistence, and lateral movement. 
Mastering these logs is essential for threat hunting, detection engineering, and CTI validation. 
Later chapters build on these fundamentals to create full hunting playbooks, detection logic, and advanced patterns 
that reflect real-world analyst workflows.
